%! Introduction to Polynomial Functions — Unit 2, Lesson 2.0 — Homework
\documentclass[10pt]{article}
\usepackage{precalculus-article}
\usepackage{precalculus-boxes}
\newcommand{\TallMath}[1]{$\displaystyle #1\rule[-1.4em]{0pt}{3.2em}$}

\begin{document}

\pageheader{Unit 2, Lesson 2.0}{Homework}

\vspace{0.12in}

\begin{enumerate}[itemsep=10pt]

\item Name the parts of each polynomial.

\smallskip
\renewcommand{\arraystretch}{1.6}
\begin{tabular}{|p{0.26\linewidth}|c|c|c|c|}
\hline
\textbf{Polynomial} & \textbf{Degree} & \textbf{Leading term} &
\textbf{Leading coef.} & \textbf{Constant term} \\
\hline
$p(x) = 3x^5 - x^2 + 8$ & \blank{1.2cm} & \blank{1.8cm} & \blank{1.4cm} & \blank{1.4cm} \\
\hline
$q(x) = -x^2 + 4x$ & \blank{1.2cm} & \blank{1.8cm} & \blank{1.4cm} & \blank{1.4cm} \\
\hline
\end{tabular}

\item Is each expression a polynomial? Circle one and give a short reason.

\begin{enumerate}[label=(\alph*), itemsep=6pt]
  \item $f(x) = 2x^3 - 9x + 4$ \quad \textbf{Yes} / \textbf{No}

        \vspace{0.05in}
        \writeline

  \item $g(x) = \dfrac{7}{x^2} + x$ \quad \textbf{Yes} / \textbf{No}

        \vspace{0.05in}
        \writeline

  \item $h(x) = 4\sqrt{x} + 1$ \quad \textbf{Yes} / \textbf{No}

        \vspace{0.05in}
        \writeline

  \item $k(x) = \dfrac{x^4}{3} - 2x$ \quad \textbf{Yes} / \textbf{No}

        \vspace{0.05in}
        \writeline
\end{enumerate}

\item Write each polynomial in standard form and name its degree family.

\smallskip
\renewcommand{\arraystretch}{1.6}
\begin{tabular}{|p{0.24\linewidth}|p{0.30\linewidth}|p{0.24\linewidth}|}
\hline
\textbf{As given} & \textbf{Standard form} & \textbf{Degree family} \\
\hline
$6 - 4x$ & \blank{3.2cm} & \blank{2.6cm} \\
\hline
$5x^2 + x^3 - 1$ & \blank{3.2cm} & \blank{2.6cm} \\
\hline
$-3 + 2x^2$ & \blank{3.2cm} & \blank{2.6cm} \\
\hline
\end{tabular}

\item Two polynomial graphs are shown. Both are smooth and unbroken.

\begin{center}
\begin{tabular}{cc}
\begin{tikzpicture}[x=0.50cm, y=0.34cm]
  \draw[->, thick] (-3.2,0) -- (3.2,0);
  \draw[->, thick] (0,-3.2) -- (0,3.4);
  \draw[very thick, plum, domain=-2.6:2.6, samples=80, smooth, variable=\x]
    plot ({\x}, {-0.5*\x*\x*\x + 1.5*\x});
\end{tikzpicture}
&
\begin{tikzpicture}[x=0.50cm, y=0.34cm]
  \draw[->, thick] (-3.2,0) -- (3.2,0);
  \draw[->, thick] (0,-3.2) -- (0,3.4);
  \draw[very thick, plum, domain=-2.6:2.6, samples=80, smooth, variable=\x]
    plot ({\x}, {0.5*\x*\x - 2});
\end{tikzpicture}
\\[2pt]
\textbf{Graph P} & \textbf{Graph Q}
\end{tabular}
\end{center}

\begin{enumerate}[label=(\alph*), itemsep=4pt]
  \item Graph P has \blank{1.4cm} turning points, so its polynomial has degree at
        least \blank{1.4cm}.
  \item Graph Q has \blank{1.4cm} turning points, so its polynomial has degree at
        least \blank{1.4cm}.
  \item Name one feature that neither graph has, because no polynomial graph ever
        has it. \blank{4cm}
\end{enumerate}

\item Which of the following is \textbf{NOT} a polynomial function?

\begin{enumerate}[label=\Alph*., itemsep=3pt]
  \item $p(x) = x^4 - 3x^2 + 7$
  \item $p(x) = 5$
  \item $p(x) = 2x^{-3} + x$
  \item $p(x) = -x^3 + \dfrac{x}{2}$
\end{enumerate}

\end{enumerate}

\vspace{0.1in}

\boxguard
\begin{extensionbox}
A polynomial graph is smooth and unbroken and has exactly 4 turning points. What is
the smallest degree its polynomial could have? Could its degree be 7 instead?
Explain both answers.

\vspace{0.05in}
\writelines{2}
\end{extensionbox}

\vspace{0.1in}

\boxguard
\begin{notesbox}{Preview of Lesson 2.1}
Today you learned to recognize the polynomial family and name its parts. Next
lesson, \textbf{Quadratic Functions Revisited}, we zoom in on the degree-2 members
you already know from Algebra 2 --- vertex, axis of symmetry, and the two forms a
quadratic can be written in --- and use them as the model for everything the rest
of the unit does with higher degrees.
\end{notesbox}

\end{document}
